\documentclass{article}
\usepackage[utf8]{inputenc}
\usepackage{amsmath}
\usepackage{geometry}
\geometry{margin=2cm}
\begin{document}

\section*{Questão 133 — ENEM 2024 — Ciências da Natureza}

O Cerrado e a Amazônia abrigam grande número de serpentes popularmente conhecidas como cobras-corais. Na Amazônia predominam as corais-verdadeiras, que são peçonhentas, enquanto no Cerrado prevalecem as falsas-corais, que não possuem peçonha. Essas espécies apresentam um padrão de coloração muito semelhante.

Nas fotografias, são apresentados exemplos dessas serpentes: uma coral-verdadeira e uma falsa-coral.

\begin{center}
  \textit{(Imagens indicadas no enunciado mostram a semelhança na coloração entre exemplares das espécies \textit{Micrurus decoratus} e \textit{Erythrolamprus aesculapii})}
\end{center}

\section*{Interpretação e Estratégia}

A questão demanda distinguir qual vantagem essa similaridade proporciona às falsas-corais. Essa análise requer o conhecimento do mimetismo Batesiano, fenômeno biológico em que uma espécie inofensiva adota a aparência de outra venenosa para se proteger.

\begin{itemize}
  \item Corais-verdadeiras são peçonhentas.
  \item Falsas-corais não possuem peçonha.
  \item Ambas apresentam coloração semelhante.
  \item Predadores evitam animais com padrões associados a perigo.
  \item Assim, as falsas-corais são menos predadas ao se parecerem com as verdadeiras.
\end{itemize}

\section*{Mini Revisão e Conteúdo}

\begin{tabular}{|p{4cm}|p{6cm}|p{5cm}|}
  \hline
  \textbf{Conceito} & \textbf{Ideia Principal} & \textbf{Aplicação ao Contexto} \\
  \hline
  Mimetismo Batesiano & Espécies inofensivas imitam espécies perigosas & Falsas-corais imitam corais-verdadeiras para proteção \\
  \hline
  Predação & Predadores evitam padrões de perigo visuais & Falsas-corais se beneficiam da aparência venenosa \\
  \hline
  Sobrevivência & Aparência influencia interações ecológicas & Coloração proporciona vantagem defensiva \\
  \hline
  Erros Comuns & Mimetismo não auxilia caça nem reprodução & Alternativas incorretas refletem essas confusões \\
  \hline
\end{tabular}

\bigskip

\textbf{Fundamentação teórica:}

O mimetismo é uma estratégia evolutiva que confere maior chance de sobrevivência à espécie mimetizadora, que apresenta aparência semelhante a uma espécie com algum mecanismo de defesa, como veneno. No mimetismo Batesiano, a espécie inofensiva se protege imitando a perigosa, reduzindo ataques de predadores através dessa confusão visual.

\section*{Resolução e ``Pulo do Gato''}

As corais-verdadeiras são peçonhentas, mas as falsas-corais não. A coloração semelhante funciona como uma proteção para as falsas-corais, pois predadores aprendem a evitar o padrão característico das serpentes venenosas.

Logo, a vantagem da similaridade para as falsas-corais é a redução das chances de serem atacadas por predadores.

\medskip

\textbf{Pulo do gato:} Espécies inofensivas podem garantir sua proteção imitando a aparência de espécies perigosas — isso é mimetismo Batesiano.

\section*{Análise das Alternativas Incorretas}

\begin{itemize}
  \item \textbf{A}: Errado — a semelhança não facilita a caça das falsas-corais.
  \item \textbf{B}: Errado — a coloração não reduz competição entre as espécies.
  \item \textbf{C}: Errado — aparência parecida não indica possibilidade de híbridos.
  \item \textbf{E}: Errado — semelhança não promove encontro reprodutivo entre espécies distintas.
\end{itemize}

\section*{Código da Questão}

CN\_ECO\_MIMETISMO\_001

% Referências das imagens do enunciado:
% SILVA, L. C.; COTTA, G. A.; RESENDE, F. C. "Cobra-coral: aplicativo educativo para reconhecimento das cobras-corais do estado de Minas Gerais, Brasil." Herpetologia Brasileira, n. 1, 2021.

\end{document}